% Options for packages loaded elsewhere
% Options for packages loaded elsewhere
\PassOptionsToPackage{unicode}{hyperref}
\PassOptionsToPackage{hyphens}{url}
\PassOptionsToPackage{dvipsnames,svgnames,x11names}{xcolor}
%
\documentclass[
]{agujournal2019}
\usepackage{xcolor}
\usepackage{amsmath,amssymb}
\setcounter{secnumdepth}{5}
\usepackage{iftex}
\ifPDFTeX
  \usepackage[T1]{fontenc}
  \usepackage[utf8]{inputenc}
  \usepackage{textcomp} % provide euro and other symbols
\else % if luatex or xetex
  \usepackage{unicode-math} % this also loads fontspec
  \defaultfontfeatures{Scale=MatchLowercase}
  \defaultfontfeatures[\rmfamily]{Ligatures=TeX,Scale=1}
\fi
\usepackage{lmodern}
\ifPDFTeX\else
  % xetex/luatex font selection
\fi
% Use upquote if available, for straight quotes in verbatim environments
\IfFileExists{upquote.sty}{\usepackage{upquote}}{}
\IfFileExists{microtype.sty}{% use microtype if available
  \usepackage[]{microtype}
  \UseMicrotypeSet[protrusion]{basicmath} % disable protrusion for tt fonts
}{}
\makeatletter
\@ifundefined{KOMAClassName}{% if non-KOMA class
  \IfFileExists{parskip.sty}{%
    \usepackage{parskip}
  }{% else
    \setlength{\parindent}{0pt}
    \setlength{\parskip}{6pt plus 2pt minus 1pt}}
}{% if KOMA class
  \KOMAoptions{parskip=half}}
\makeatother
% Make \paragraph and \subparagraph free-standing
\makeatletter
\ifx\paragraph\undefined\else
  \let\oldparagraph\paragraph
  \renewcommand{\paragraph}{
    \@ifstar
      \xxxParagraphStar
      \xxxParagraphNoStar
  }
  \newcommand{\xxxParagraphStar}[1]{\oldparagraph*{#1}\mbox{}}
  \newcommand{\xxxParagraphNoStar}[1]{\oldparagraph{#1}\mbox{}}
\fi
\ifx\subparagraph\undefined\else
  \let\oldsubparagraph\subparagraph
  \renewcommand{\subparagraph}{
    \@ifstar
      \xxxSubParagraphStar
      \xxxSubParagraphNoStar
  }
  \newcommand{\xxxSubParagraphStar}[1]{\oldsubparagraph*{#1}\mbox{}}
  \newcommand{\xxxSubParagraphNoStar}[1]{\oldsubparagraph{#1}\mbox{}}
\fi
\makeatother


\usepackage{longtable,booktabs,array}
\usepackage{calc} % for calculating minipage widths
% Correct order of tables after \paragraph or \subparagraph
\usepackage{etoolbox}
\makeatletter
\patchcmd\longtable{\par}{\if@noskipsec\mbox{}\fi\par}{}{}
\makeatother
% Allow footnotes in longtable head/foot
\IfFileExists{footnotehyper.sty}{\usepackage{footnotehyper}}{\usepackage{footnote}}
\makesavenoteenv{longtable}
\usepackage{graphicx}
\makeatletter
\newsavebox\pandoc@box
\newcommand*\pandocbounded[1]{% scales image to fit in text height/width
  \sbox\pandoc@box{#1}%
  \Gscale@div\@tempa{\textheight}{\dimexpr\ht\pandoc@box+\dp\pandoc@box\relax}%
  \Gscale@div\@tempb{\linewidth}{\wd\pandoc@box}%
  \ifdim\@tempb\p@<\@tempa\p@\let\@tempa\@tempb\fi% select the smaller of both
  \ifdim\@tempa\p@<\p@\scalebox{\@tempa}{\usebox\pandoc@box}%
  \else\usebox{\pandoc@box}%
  \fi%
}
% Set default figure placement to htbp
\def\fps@figure{htbp}
\makeatother


% definitions for citeproc citations
\NewDocumentCommand\citeproctext{}{}
\NewDocumentCommand\citeproc{mm}{%
  \begingroup\def\citeproctext{#2}\cite{#1}\endgroup}
\makeatletter
 % allow citations to break across lines
 \let\@cite@ofmt\@firstofone
 % avoid brackets around text for \cite:
 \def\@biblabel#1{}
 \def\@cite#1#2{{#1\if@tempswa , #2\fi}}
\makeatother
\newlength{\cslhangindent}
\setlength{\cslhangindent}{1.5em}
\newlength{\csllabelwidth}
\setlength{\csllabelwidth}{3em}
\newenvironment{CSLReferences}[2] % #1 hanging-indent, #2 entry-spacing
 {\begin{list}{}{%
  \setlength{\itemindent}{0pt}
  \setlength{\leftmargin}{0pt}
  \setlength{\parsep}{0pt}
  % turn on hanging indent if param 1 is 1
  \ifodd #1
   \setlength{\leftmargin}{\cslhangindent}
   \setlength{\itemindent}{-1\cslhangindent}
  \fi
  % set entry spacing
  \setlength{\itemsep}{#2\baselineskip}}}
 {\end{list}}
\usepackage{calc}
\newcommand{\CSLBlock}[1]{\hfill\break\parbox[t]{\linewidth}{\strut\ignorespaces#1\strut}}
\newcommand{\CSLLeftMargin}[1]{\parbox[t]{\csllabelwidth}{\strut#1\strut}}
\newcommand{\CSLRightInline}[1]{\parbox[t]{\linewidth - \csllabelwidth}{\strut#1\strut}}
\newcommand{\CSLIndent}[1]{\hspace{\cslhangindent}#1}



\setlength{\emergencystretch}{3em} % prevent overfull lines

\providecommand{\tightlist}{%
  \setlength{\itemsep}{0pt}\setlength{\parskip}{0pt}}



 


\usepackage{url} %this package should fix any errors with URLs in refs.
\usepackage{lineno}
\usepackage[inline]{trackchanges} %for better track changes. finalnew option will compile document with changes incorporated.
\usepackage{soul}
\linenumbers
\makeatletter
\@ifpackageloaded{caption}{}{\usepackage{caption}}
\AtBeginDocument{%
\ifdefined\contentsname
  \renewcommand*\contentsname{Table of contents}
\else
  \newcommand\contentsname{Table of contents}
\fi
\ifdefined\listfigurename
  \renewcommand*\listfigurename{List of Figures}
\else
  \newcommand\listfigurename{List of Figures}
\fi
\ifdefined\listtablename
  \renewcommand*\listtablename{List of Tables}
\else
  \newcommand\listtablename{List of Tables}
\fi
\ifdefined\figurename
  \renewcommand*\figurename{Figure}
\else
  \newcommand\figurename{Figure}
\fi
\ifdefined\tablename
  \renewcommand*\tablename{Table}
\else
  \newcommand\tablename{Table}
\fi
}
\@ifpackageloaded{float}{}{\usepackage{float}}
\floatstyle{ruled}
\@ifundefined{c@chapter}{\newfloat{codelisting}{h}{lop}}{\newfloat{codelisting}{h}{lop}[chapter]}
\floatname{codelisting}{Listing}
\newcommand*\listoflistings{\listof{codelisting}{List of Listings}}
\makeatother
\makeatletter
\makeatother
\makeatletter
\@ifpackageloaded{caption}{}{\usepackage{caption}}
\@ifpackageloaded{subcaption}{}{\usepackage{subcaption}}
\makeatother
\usepackage{bookmark}
\IfFileExists{xurl.sty}{\usepackage{xurl}}{} % add URL line breaks if available
\urlstyle{same}
\hypersetup{
  pdftitle={Using rapid temperature falls to estimate future strong cold front frequency in CMIP6 climate projections},
  pdfauthor={Paola B Corrales; Martin Singh; Michael J. Reeder},
  colorlinks=true,
  linkcolor={blue},
  filecolor={Maroon},
  citecolor={Blue},
  urlcolor={Blue},
  pdfcreator={LaTeX via pandoc}}


\journalname{Geophysical Research Letters}

\begin{document}
\title{Using rapid temperature falls to estimate future strong cold
front frequency in CMIP6 climate projections}

\authors{Paola B Corrales\affil{1,2,3,4}, Martin
Singh\affil{1,2,4}, Michael J. Reeder\affil{1,2}}
\affiliation{1}{School of Earth, Atmosphere and Environment, Monash
University, Melbourne, Australia}\affiliation{2}{ARC Centre of
Excellence for Climate
Extremes, Melbourne, Australia}\affiliation{3}{Research School of Earth
Sciences, Australian National
University, Canberra, Australia}\affiliation{4}{ARC Centre of Excellence
for the Weather of the 21st Century, Melbourne, Australia}
\correspondingauthor{Paola B Corrales}{paola.corrales@monash.edu}

\begin{keypoints}
\item Day-to-day falls in maximum temperature are used to identify cold
fronts in CMIP6 models using the historical and SSP5-8.5
simulations \item In the NH the frequency of strong cold fronts
decreases, while in the SH they shift poleward over the oceans and
increase over land \item Cold front frequency changes correlated
positively with the meridional surface temperature gradient globally and
with storm tracks in the NH 
\end{keypoints}

\begin{abstract}
We use day-to-day falls in daily-maximum surface air temperature to
assess future changes in the frequency of strong cold fronts in a set of
CMIP6 climate models under the SSP5-8.5 scenario. We find a general
decrease in cold front frequency across the Northern Hemisphere,
particularly over the oceans, and in the Southen Hemisphere a poleward
shift in the region of highest frontal frequency, with increased
frequency over midlatitude land regions. The frequency of strong
temperature falls correlates positively with the meridional temperature
gradient and transient eddy kinetic energy, implicating large-scale
baroclinicity and changes in the storm track influence the pattern of
strong cold fronts. Over Southern Hemisphere land areas, this
relationship weakens, suggesting the influence of regional thermodynamic
processes.
\end{abstract}

\section*{Plain Language Summary}
Cold fronts, are defined by large temperature falls and sometimes
accompanied by rain and strong wind. They play a key role in the
day-to-day weather in midlatitudes. In this study, we use a simple
method based on changes in maximun temperature from one day to the next
to identify strong cold fronts in climate simulations. In the future
climate projections from ten global models, strong cold fronts are less
frequent in much of the Northern Hemisphere, especially over the oceans.
In the Southern Hemisphere, however, strong cold fronts in the
projections shift closer to the South Pole and may become more frequent
over land areas like southern Australia and southern Africa. These
changes are linked to shifts in the temperature contrast between the
equator and the poles and the paths taken by storm systems.




\section{Introduction}\label{introduction}

Cold fronts are an important feature of weather in extratropical regions
and their associated circulation contribute a substantial fraction of
the precipitation in these regions (Catto \& Pfahl, 2013; Konstali et
al., 2024). Cold fronts are at times linked to extreme events, such as
heavy precipitation (Catto \& Pfahl, 2013) and severe wildfires in
southern Australia (Reeder et al., 2015) and other regions.
Understanding how cold fronts, and especially strong cold fronts, change
in a warming climate is therefore of importance for understanding the
changing likelihood of extremes and fire weather.

Several computational methods have been developed to identify cold
fronts using different characteristics. Some of these methods look for
spatial gradients in temperature and humidity at different atmospheric
levels (Berry et al., 2011), others identify gradients in the wind field
alone (Simmonds et al., 2012) or a combination of dynamic and
thermodynamic variables (Thomas \& Schultz, 2019). These methods have
been developed for the data on a regular grid (e.g.~outputs from
numerical models or reanalysis) and their results depend on setting
different thresholds, which may need to be calibrated for different grid
resolutions.

Here, we take the simplest approach, first introduced by Reeder et al.
(2015). We measure the occurrence of cold fronts using temporal changes
in temperature at a fixed location. Specifically, we identify the
passage of cold fronts with large day-to-day falls in near-surface
daily-maximum temperature. This method has been shown to be suitable for
the detection of strong cold fronts in southeastern Australia (Cai et
al., 2022; Reeder et al., 2015), and can successfully identify cold
fronts in data with different resolutions, including both station data
and reanalyses (King et al., 2024). King et al. (2024) developed a
climatology of cold-front frequency in global climate models from phase
6 of the Coupled Model Intercomparison Project (CMIP6) which produced
results comparable to other, more complex, cold-front detection
algorithms. While the models tend to overestimate front frequency in
polar regions and underestimate the occurrence of large temperature
falls over midlatitude land regions compared with ERA5, they accurately
represent the patterns of front frequency equatorward of about 60°.

As global temperatures continue to rise, understanding how this warming
will affect the occurrence of cold fronts is crucial for projecting
changes in extreme weather and fire risk. One of the few studies to
address this question is Catto et al. (2014), who identified fronts
using an objective detection method based on temperature and humidity
fields (Berry et al., 2011). In this framework, fronts are represented
as line features, marking sharp gradients in thermodynamic variables.
Using this approach in 18 CMIP5 models, they compared historical
simulations (1980--2005) with projections under the Representative
Concentration Pathway 8.5 scenario (2080--2100). They found an overall
decrease in front frequency in the Northern Hemisphere, with a poleward
shift in the latitude of maximum frequency and a significant decrease at
high latitudes. The Southern Hemisphere also showed a poleward shift in
the frequency maximum.

In contrast, Konstali et al. (2024) identified fronts in CESM2-LE
simulations using the gradient of equivalent potential temperature at
850 hPa. In this approach, frontal systems are represented as areas of
strong gradients rather than discrete lines. Using this methodology,
they found an increase in the frequency of fronts, which they suggest
may be associated with fronts becoming stronger or more elongated in
both hemispheres.

Motivated by the contrasting results of previous studies, we re-examine
how the distribution of cold fronts change in projections from the CMIP6
ensemble, but we consider a simpler method of estimating changes in
front frequency and intensity with cooling based on day-to-day falls in
maximum surface air temperature. This approach offers several
advantages. First, the method's simplicity and weak
resolution-dependence (as shown by King et al., 2024) allows us to
compare it with earlier studies using other methodologies. Second,
focusing on local temperature falls provides a complementary perspective
to Lagrangian approaches, as day-to-day temperature falls are the most
direct impact of cold fronts on local weather.

Consistent with Catto et al. (2014), we find a decrease in cold front
frequency in the Northern Hemisphere and a poleward shift in the
Southern Hemisphere. Seeking to develop a physical explanation for these
changes, we first study the response of the meridional temperature
gradient and its relationship with the change in front frequency. We
also explore the relationship between changes in front frequency and
transient eddy kinetic energy, as it is directly related to the storm
tracks. Several studies have shown a poleward shift on the storm tracks
in future climate change scenarios, particularly in the Southern
Hemisphere (Chemke, 2022; O'Gorman, 2010; Tamarin \& Kaspi, 2017), and
these changes could also influence the location and intensity of future
cold fronts.

\section{Data and Methods}\label{data-and-methods}

We analyse 10 models from phase 6 of the Coupled Model Intercomparison
Project (CMIP6, Eyring et al. (2016)), including the Historical
experiment up to 2014 and the Shared Socioeconomic Pathways 5 (SSP5-8.5)
from the Scenario Model Intercomparison Project (ScenarioMIP, O'Neill et
al. (2016)) from 2015 to 2100, as described in Table~\ref{tbl-models}.
We selected these models based on the data available for the historical
and scenario experiments stored on the Australian National Computational
Infrastructure (NCI) replica data store.

\begin{longtable}[]{@{}
  >{\raggedright\arraybackslash}p{(\linewidth - 4\tabcolsep) * \real{0.6273}}
  >{\raggedright\arraybackslash}p{(\linewidth - 4\tabcolsep) * \real{0.1818}}
  >{\raggedright\arraybackslash}p{(\linewidth - 4\tabcolsep) * \real{0.1909}}@{}}
\caption{List of the coupled model intercomparison project phase 6
(CMIP6) climate model used in this
analysis.}\label{tbl-models}\tabularnewline
\toprule\noalign{}
\begin{minipage}[b]{\linewidth}\raggedright
Model
\end{minipage} & \begin{minipage}[b]{\linewidth}\raggedright
Institution
\end{minipage} & \begin{minipage}[b]{\linewidth}\raggedright
Ensamble member used
\end{minipage} \\
\midrule\noalign{}
\endfirsthead
\toprule\noalign{}
\begin{minipage}[b]{\linewidth}\raggedright
Model
\end{minipage} & \begin{minipage}[b]{\linewidth}\raggedright
Institution
\end{minipage} & \begin{minipage}[b]{\linewidth}\raggedright
Ensamble member used
\end{minipage} \\
\midrule\noalign{}
\endhead
\bottomrule\noalign{}
\endlastfoot
CMCC-ESM2 & & \\
(Lovato et al., 2021a, 2021b) & CMCC & r1i1p1f1 \\
EC-Earth3 & & \\
(Consortium (EC-Earth), 2019a, 2019b) & EC-Earth-Consortium &
r1i1p1f1 \\
EC-Earth3-CC & & \\
(Consortium (EC-Earth), 2020a, 2021) & EC-Earth-Consortium & r1i1p1f1 \\
EC-Earth3-Veg & & \\
(Consortium (EC-Earth), 2019c, 2019d) & EC-Earth-Consortium &
r1i1p1f1 \\
EC-Earth3-Veg-LR & & \\
(Consortium (EC-Earth), 2020b, 2023) & EC-Earth-Consortium & r1i1p1f1 \\
GFDL-CM4 & & \\
(Guo, John, Blanton, McHugh, Nikonov, Radhakrishnan, Rand, Zadeh,
Balaji, Durachta, Dupuis, Menzel, Robinson, Underwood, Vahlenkamp,
Bushuk, et al., 2018; 2018) & NOAA-GFDL & r1i1p1f1 \\
INM-CM4-8 & & \\
Volodin et al. (2019b) & INM & r1i1p1f1 \\
INM-CM5-0 & & \\
(Volodin et al., 2019c, 2019d) & INM & r1i1p1f1 \\
MPI-ESM1-2-HR & & \\
(Jungclaus et al., 2019; Schupfner et al., 2019) & DKRZ & r1i1p1f1 \\
NorESM2-MM & & \\
(Bentsen et al., 2019a, 2019b) & NCC & r1i1p1f1 \\
\end{longtable}

\subsection{Front Identification}\label{front-identification}

We identify falls in near-surface (2 meter) air temperature as a proxy
for frontal passages by calculating the change in daily maximum
temperature from one day to the next (\(\Delta T\)). \(\Delta T\) is
calculated for all grid points and days in the study period. King et al.
(2024) defined a passage of a strong cold front as
\(\Delta T \le -10^{\circ}C\). That is, a cold front passage is
identified on day \(i\) if the maximum temperature on day \(i+1\) is at
least 10 deg C lower than that of day \(i\). While this value is
reasonable for detecting strong falls in temperature over land, it may
underestimate the frequency of fronts over the ocean due to the weaker
temperature variability. For this reason, in this study, we use the
2.5\(^{th}\) percentile of \(\Delta T\) as a threshold for each grid
point to identify strong falls in temperature, hereafter
\(\Delta T_{2.5} = \Delta T \le 2.5^{th}\) percentile. This percentile
is chosen because the mean of the 2.5\(^{th}\) percentile over all
midlatitude land areas (between 25\(^{\circ}\) to 55\(^{\circ}\)) is
close to \(-10^{\circ}C\), the threshold used by King et al. (2024). The
percentiles are calculated over the historical period between 1979 and
2014 for each model at each grid point. Additionally, the percentile
values are also calculated over the same period for
December-January-February (DJF) and June-July-August (JJA) for the
seasonal analysis (see Supplemental Information).

It is important to note that by using this threshold we detect frontal
passages in exactly 2.5\% of days. In regions of regular frontal passage
the 2.5\(^{th}\) percentile of \(\Delta T\) generally corresponds to
strong fronts, and our identification method produces reasonable
results. However, outside these regions (e.g., in the tropics) we
identify temperature falls that can be associated with midlatitude
front-like features but also other processes.

\subsection{Time-mean Meridional Temperature
Gradient}\label{time-mean-meridional-temperature-gradient}

We calculated the time-mean meridional temperature gradient
(\(\partial \bar{T}/ \partial y\)) at each grid point using central
finite differences for the temperature at 850 hPa. We then area-averaged
\(\partial \bar{T}/ \partial y\) over different regions, e.g., between
25\(^{\circ}\) and 65\(^{\circ}\) latitude for each hemisphere. This
process produces the same results as if calculating the temperature
difference between two latitudes but has the advantage of allowing the
calculation of \(\partial \bar{T}/ \partial y\) for regions over land
and sea independently.

\subsection{Storm Tracks Estimation}\label{sec-storm-track}

To estimate the storm track location we follow previous work (Chemke,
2022; O'Gorman, 2010) and use the high-pass filtered transient eddy
kinetic energy (EKE) in the upper troposphere. For each model, we
calculated the EKE as \(EKE = 0.5(u'^2 + v'^2)\), where \(u'\) and
\(v'\) are obtained using a Fast Fourier Transform high-pass filter
applied to daily horizontal wind at 250 hPa with a 6-day cutoff.

\section{Results}\label{results}

Figure~\ref{fig-trend} a) shows the change in the frequency of
\(\Delta T_{2.5}\) between the yearly averaged \(\Delta T_{2.5}\)
frequency for last 20 years of the 21\(^{st}\) century (SSP5-8.5
experiment) and the last 20 years of the 20\(^{th}\) century (historical
experiment). Figure~\ref{fig-trend} b and c) show the zonal mean of
\(\Delta T_{2.5}\) frequency for each latitude over sea and over land
respectively.

Figure~\ref{fig-trend} d) shows the \(\Delta T\) values that define the
2.5\(^{th}\) percentile. We found that the spatial pattern correlates
with the climatology of cold fronts from Berry et al. (2011) (their
Figure 2a) and King et al. (2024) (their Figure 6). This consistency
supports our proposition that temperature falls identified here are
directly related to cold front passages.

In terms of projected changes, in the Northern Hemisphere over land,
there is a robust decrease in \(\Delta T_{2.5}\), particularly north of
40\(^{\circ}\) latitude, except in western and central Europe where the
multi-model mean shows an increase but the signal is not robust across
the ensemble. A similar pattern of changes in front frequency was also
found by Catto et al. (2014) using CMIP5 and a different methodology to
identify fronts, reinforcing the reliability of the methodology used in
this analysis. Over the oceans, front frequency also tends to decrease
over the North Atlantic, which shows a multi-model mean increase in
front frequency west of continental Europe. This local increase may be
interpreted as a shift of front frequency toward the pole in this
sector, although it does not appear consistently across the ensemble.

In the Southern Hemisphere over land, the change is generally positive
with an increase of 5 strong temperature falls per year in some areas.
Other authors have also found an increase in the frequency of fronts in
the south and east regions of Africa and Australia (Catto et al., 2014;
Konstali et al., 2024).

Changes in \(\Delta T_{2.5}\) frequency over the Southern Hemisphere
ocean regions are less robust across the ensemble, but the ensemble mean
shows a tripolar pattern, with increases in a band centered around
40\(^{\circ}\)-50\(^{\circ}\) S (depending on the longitude) and
decreases poleward and equatorward. This pattern indicates that the
areas of higher front frequency shift poleward in the future scenario.
In the Indian ocean, the pattern in more complicated, with a second band
of increased frequency in the subtropics is most pronounced in the
Agulas region.

The analysis of the change in frequency of \(\Delta T_{2.5}\) by season
in DJF and JJA (Figure S1 a) shows similar spatial patterns and the same
sign of the change when compared with the change per year in
Figure~\ref{fig-trend} a). The main exception is western Europe, where
we find an increase in front frequency during winter but a decrease in
summer.

\begin{figure}

\centering{

\pandocbounded{\includegraphics[keepaspectratio]{fig-trend-1.pdf}}

}

\caption{\label{fig-trend}a) Multimodel mean change in the frequency of
\(\Delta T_{2.5}\) (days/year, SSP5-8.5 2080-2100 minus Historical
1980-2000) and zonal mean change over b) land and c) sea, each grey line
represents a model listed in Table 1 and a black line represents the
multimodel mean. In a) dotted areas show where at least 90\% of the
models agree in the sign of change. d) \(\Delta T\) values associated
with the 2.5\(^{th}\) percentile and zonal mean over e) land and f)
sea.}

\end{figure}%

Figure~\ref{fig-serie} shows ensemble-mean time series of the zonal-mean
frequency of strong temperature falls for different latitude bands in
each hemisphere. Consistent with the map in Figure~\ref{fig-trend}, the
trends in frequency are negative over ocean in the Northern Hemisphere,
but positive over land equatorward of 50\(^{\circ}\) N. Over the
Southern Hemisphere, trends are generally positive over land, with
changes over ocean generally consistent with a poleward shift of fronts.

While there is considerable variation in these trends between models
(see boxplots on right of Figure~\ref{fig-serie}), particularly at more
poleward locations, for most of the zonal mean trends, there is
consistency in the sign of the trend across models.

\begin{figure}

\centering{

\pandocbounded{\includegraphics[keepaspectratio]{fig-serie-1.pdf}}

}

\caption{\label{fig-serie}Multimodel yearly mean frequency of
\(\Delta T_{2.5}\) for the historical period (1900-2014) and SSP5-8.5
(2015-2100) over land and sea for different 2° wide latitude bands
centered at 30, 40 and 50°. The grey line represents the climatological
frecuency for 1979-2014, equal to 9.125 \(\Delta T_{2.5}\)/year. On the
right of each panel box plots of the intermodel distribution is shown
for the period 2080-2100. In the boxplots the lower and upper hinges
correspond to the first and third quartiles and the upper (lower)
whisker extends to the largest (smallest) value no further than 1.5 *
inter-quartile range. Outside the whiskers dots represents outliers.}

\end{figure}%

To explore the distribution of the change in temperature falls of
different intensity, Figure~\ref{fig-distribution} shows the change in
the frequency of \(\Delta T\) using different percentiles, from
10\(^{th}\) to 0.1\(^{th}\), as thresholds to identify the falls in
temperature. The percentile values are again calculated from the
historical period between 1979 and 2014 for each grid point and model as
in Figure~\ref{fig-trend}.

Figure~\ref{fig-distribution} shows that, in most cases, the sign of the
change is the same across all percentile thresholds, meaning that there
is a projected increase or decrease in the frequency of \(\Delta T\)
regardless of the intensity of the temperature fall. The exceptions are
the Southern Hemisphere over sea and the Northern Hemisphere over land
in the 40 \(^{\circ}\) latitude band, where there is a modest increase
in strong temperature falls but a small decrease in weaker falls.

Moreover the change in the frequency of \(\Delta T\) is greater at lower
percentile thresholds, which identify more extreme temperature falls.
For example, in the Northern Hemisphere over the sea, the 0.1\(^{th}\)
percentile threshold shows that there will be less than half as many
falls in temperature in the future over 50\(^{\circ}\) N, with the
opposite expected over land in the Southern Hemisphere. However, lower
percentile thresholds are associated with rarer events, meaning fewer
occurrences per year. These results suggest that regions projected to
experience more temperature falls in the future, such as land areas in
midlatitudes in the Southern Hemisphere, will also experience larger
\(\Delta T\), whereas regions with a decreasing frequency of falls will
see smaller falls.

\begin{figure}

\centering{

\pandocbounded{\includegraphics[keepaspectratio]{fig-distribution-1.pdf}}

}

\caption{\label{fig-distribution}Fractional change in the frequency of
strong \(\Delta T\) averaged between the years 2080-2100 under the
SSP5-8.5 scenario compared to that averaged between 1980-2000 in the
historical scenario, for different thresholds defined by their
percentile value for each grid point and model calculated over the
historical period (1979-2014).}

\end{figure}%

\subsection{Mechanisms}\label{mechanisms}

We now investigate the mechanisms driving the changes in
\(\Delta T_{2.5}\) frequency described above. We take a simple
perspective on midlatitude temperature variability based on advection as
a key driver (Schneider et al., 2015). According to this view, changes
in midlatitude temperature variability are closely linked to changes in
large-scale temperature gradients and to the strength of mixing by
geostrophic turbulence. Motivated by this perspective, we explore
relationships between changes in front frequency and changes in the
large-scale meridional temperature gradient
(\(\partial \bar{T} / \partial y\)) and EKE within the midlatitude storm
tracks.

Figure~\ref{fig-eke} a) shows that, in both hemispheres, there is a
strong correlation between increases in front frequency and changes in
the large-scale meridional surface temperature gradient
\(\partial \bar{T} / \partial y\) across models. Here,
\(\partial \bar{T} / \partial y\) is taken as the time and zonal-mean
meridional temperature gradient averaged between 25\(^{\circ}\) and
65\(^{\circ}\) in each hemisphere over either land or ocean. As
previously noted by Catto et al. (2014), reductions in front frequency
in the Northern Hemisphere are associated with a decrease in the the
temperature gradient near the surface, while opposite changes in surface
temperature gradients in the Southern Hemisphere drive increases is
front frequency there. Here, we find that this relationship holds within
the model ensemble; larger increases or decreases in
\(\partial \bar{T} / \partial y\) in a given model are associated with
larger projected increases or decreases in \(\Delta T\), a relationship
that is particularly strong in the Northern Hemisphere.

There is a similar relationship across the ensemble between changes in
the frequency of \(\Delta T_{2.5}\) and changes in the energy of the
storm track. Here we estimate the storm track energy by taking the
high-pass filtered eddy kinetic energy at 250 hPa (see
Section~\ref{sec-storm-track}) and integrating over the 25\(^{\circ}\)
and 65\(^{\circ}\) latitudes in each hemisphere. Models with projected
decreases in EKE tend to project larger decreases in the frequency of
\(\Delta T_{2.5}\). The exception is the Southern Hemisphere over land
where all models project weak changes in EKE and large increases in the
frequency of \(\Delta T_{2.5}\). The reason for this difference is
unclear, but could suggest the importance of local land processes in
amplifying fronts over land in some regions of the Southern Hemisphere
(Muir \& Reeder, 2010). Moreover, most of the land affected by fronts in
the Southern Hemisphere is north of the mean storm track, and thus the
storm track has only a weak effect on temperature variability there.
Supporting this, Harvey et al. (2014) found that the storm track
response in the Southern Hemisphere during summer is associated with the
change in the low level \(\partial \bar{T} / \partial y\) mainly
poleward of 65\(^{\circ}\) S, i.e.~over the southern oceans.

In agreement with the previous literature on seasonal changes in the
midlatitude storm track (O'Gorman, 2010; Priestley \& Catto, 2022), we
find a general increase in EKE during summer and a decrease in winter in
the Northern Hemisphere. There is also a general positive change in the
Southern Hemisphere in both seasons associated with a poleward shift of
the storm tracks (Figure S2 c-d). Given these seasonal asymmetries, it
is important to note that the change in EKE averaged over the entire
year (Figure~\ref{fig-eke} b) is dominated by the winter season in each
hemisphere.

Notwithstanding the above caveats, our results suggest that changes in
the mean meridional temperature gradient and storm track intensity play
an important role in modulating changes in front frequency in a warming
climate. Indeed, differences in how these quantities change in
projections appear to account for the differences in front frequency
changes between hemispheres. However, our results do not provide strong
evidence for the relative importance of the two factors: aside from over
Southern Hemisphere, both factors correlate with the \(\Delta T_{2.5}\)
frequency across the model ensemble and both are theoretically
associated with changes in temperature variability.

\begin{figure}

\centering{

\pandocbounded{\includegraphics[keepaspectratio]{fig-eke-1.pdf}}

}

\caption{\label{fig-eke}Mean change in the mean \(\Delta T_{2.5}\)
frequency and a) mean merdional temperature gradient
\(\partial \bar{T} / \partial y\) and b) the EKE of the storm track for
SSP5-8.5 (2080-2100) compared with the historical period (1980-2000) for
each model. Values are calculated for the Northern (25\(^{\circ}\) to
65\(^{\circ}\)) and Southern (-25\(^{\circ}\) to -65\(^{\circ}\))
Hemispheres and over sea (circles) and land (triangles). Each solid line
corresponds to a linear regression for each hemisphere.}

\end{figure}%

\section{Summary and Discussion}\label{summary-and-discussion}

In this study, we use day-to-day falls in near-surface maximum
temperature to detect cold fronts, and we evaluate changes in their
frequency under future warming in a set of 10 models from CMIP6. By
using a specific percentile threshold, we ensure a consistent detection
of strong temperature falls across different regions, improving the
detection particularly over oceans where the temperature variability is
lower. While the temperature fall method may occasionally misidentify
other forms of temperature variability as fronts, previous studies have
shown that this methodology provides a simple and robust method for
characterising the frequency of fronts over midlatitude regions (King et
al., 2024).

Our results show a general decrease in the frequency of fronts across
the Northern Hemisphere, especially over oceans, and a poleward shift in
frontal activity in the Southern Hemisphere. Over land in the Southern
Hemisphere, notably in southern Australia and Africa, we find an
increase in frontal frequency and intensity. These findings align with
previous studies using more complex methods to identify fronts,
strengthening confidence in our approach.

We find a positive correlation between changes in \(\Delta T_{2.5}\)
frequency and the meridional temperature gradient and transient eddy
kinetic energy (EKE). This is consistent with a model for midlatitude
temperature variability driven primarily by advection, and suggests that
the large-scale baroclinic background state is an important control on
front frequency in a warming climate. Further work could tease apart the
individual effects of temperature gradients versus the intensity of the
storm track, noting that these are linked, since changes in the mean
thermodynamic structure affect how the storm track behaves (O'Gorman,
2010). An important exception to the above correlations is that, in
Southern Hemisphere land regions, front frequency increases despite a
decrease in the EKE of the storm track. This may reflect the fact that
storm track dynamics have less direct control over frontal systems in
this region, which could be more influenced by mesoscale processes,
land--sea contrast, and regional thermodynamic effects.

There are other factors not analyzed here that could contribute to the
projected changes in frontal activity. For example, sea surface
temperature (SST) gradients can influence the development and intensity
of atmospheric fronts through air--sea interactions. Parfitt et al.
(2017) suggested that surface sensible heat fluxes associated with
strong SST gradients, such as those found in the Gulf Stream or Kuroshio
Current regions, can amplify atmospheric fronts through diabatic
frontogenesis, although Reeder et al. (2021) showed that the effect was
only significant when the fronts were weak. In contrast, Reeder et al.
(2021) showed that strong atmospheric fronts along SST fronts are
amplified primarily by adiabatic frontogenesis (i.e.~the deformation of
the wind field).

More generally, frontal cyclogenesis is directly affected by the wind
and temperature fields. Consequently, projected changes in frontal
frequency may also reflect modifications in the large-scale circulation,
such as shifts in the background flow, changes in horizontal wind shear,
or variations in baroclinicity. Further analysis is needed to quantify
the contribution of these mechanisms to projected changes.

Finally, while this study focuses on cold fronts, in some regions warm
fronts also play an important role in the distribution of extreme
precipitation (Catto \& Pfahl, 2013). We explored extending the
methodology to warm fronts by considering day-to-day increases in
maximum temperature. However, unlike cold fronts, warm fronts are not
associated with a sufficiently distinct temperature signal to allow a
simple threshold-based identification without also including changes
associated with other local thermodynamic processes. Developing a simple
and robust method for identifying warm fronts would extend the present
framework and provide further insight into projected changes in frontal
activity under climate change.

\section{Acknowledgments}\label{acknowledgments}

This research was supported by the Australian Research Council Centre of
Excellence for Climate Extremes (CE170100023), the Australian Research
Council Centre of Excellence for the Weather of the 21st Century
(CE230100012), with the assistance of resources from the National
Computational Infrastructure (NCI Australia), an NCRIS enabled
capability supported by the Australian Government. P.C. also
acknowledges the use of the following open-source software projects that
supported data analysis, visualization, and reproducible reporting: R
(v4.3.1, R Core Team (2023)), metR (Campitelli, 2025a), rcdo
(Campitelli, 2025b), data.table (Barrett et al., 2025), ggplot2
(Wickham, 2016), patchwork (Pedersen, 2024), cowplot (Wilke, 2024),
ggside (Landis, 2024), and Quarto (Allaire et al., 2025).

\section{Conflict of Interest}\label{conflict-of-interest}

The authors declare no conflicts of interest relevant to this study.

\section{Availability Statement}\label{availability-statement}

The CMIP6 data used in the manuscript is publicly available at the Earth
System Grid Federation \url{https://esgf-node.llnl.gov/projects/cmip6/}.
A version-controlled repository of the code used to create this
analysis, including the necessary code to download the derived data has
been deposited in Zenodo (Corrales, 2026b). The derived data that
support the findings of this study are also openly available in Zenodo
(Corrales, 2026a).

\section*{References}\label{references}
\addcontentsline{toc}{section}{References}

\phantomsection\label{refs}
\begin{CSLReferences}{1}{0}
\vspace{1em}

\bibitem[\citeproctext]{ref-quarto}
Allaire, J. J., Teague, C., Scheidegger, C., Xie, Y., Dervieux, C., \&
Woodhull, G. (2025, April). {Quarto} (Version 1.7).
\url{https://doi.org/10.5281/zenodo.5960048}

\bibitem[\citeproctext]{ref-data.table}
Barrett, T., Dowle, M., Srinivasan, A., Gorecki, J., Chirico, M.,
Hocking, T., et al. (2025). Data.table: Extension of `data.frame`.

\bibitem[\citeproctext]{ref-Bentsen2019}
Bentsen, M., Oliviè, D. J. L., Seland, Ø., Toniazzo, T., Gjermundsen,
A., Graff, L. S., et al. (2019a). {IPCC DDC}: {NCC NorESM2-MM} model
output prepared for {CMIP6 CMIP} {[}Data set{]}. World Data Center for
Climate (WDCC) at DKRZ.
\url{https://doi.org/10.26050/WDCC/AR6.C6CMNCCN2}

\bibitem[\citeproctext]{ref-Bentsen2019a}
Bentsen, M., Oliviè, D. J. L., Seland, Ø., Toniazzo, T., Gjermundsen,
A., Graff, L. S., et al. (2019b). {IPCC DDC}: {NCC NorESM2-MM} model
output prepared for {CMIP6 ScenarioMIP} {[}Data set{]}. World Data
Center for Climate (WDCC) at DKRZ.
\url{https://doi.org/10.26050/WDCC/AR6.C6SPNCCN2}

\bibitem[\citeproctext]{ref-Berry2011}
Berry, G., Reeder, M. J., \& Jakob, C. (2011). A global climatology of
atmospheric fronts. \emph{Geophysical Research Letters}, \emph{38}(4).
\url{https://doi.org/10.1029/2010GL046451}

\bibitem[\citeproctext]{ref-Cai2022}
Cai, D., Abram, N. J., Sharples, J. J., \& Perkins-Kirkpatrick, S. E.
(2022). Increasing intensity and frequency of cold fronts contributed to
{Australia}'s 2019--2020 {Black Summer} fire disaster.
\emph{Environmental Research Letters}, \emph{17}(9), 094044.
\url{https://doi.org/10.1088/1748-9326/ac8e88}

\bibitem[\citeproctext]{ref-metR}
Campitelli, E. (2025a). metR: Tools for easier analysis of
meteorological fields. \url{https://doi.org/10.5281/zenodo.2593516}

\bibitem[\citeproctext]{ref-rcdo}
Campitelli, E. (2025b). Rcdo: Create CDO commands from r functions.
Retrieved from \url{https://eliocamp.github.io/rcdo/}

\bibitem[\citeproctext]{ref-Catto2013}
Catto, J. L., \& Pfahl, S. (2013). The importance of fronts for extreme
precipitation. \emph{Journal of Geophysical Research: Atmospheres},
\emph{118}(19), 10, 791--10, 801.
\url{https://doi.org/10.1002/jgrd.50852}

\bibitem[\citeproctext]{ref-Catto2014}
Catto, J. L., Nicholls, N., Jakob, C., \& Shelton, K. L. (2014).
Atmospheric fronts in current and future climates. \emph{Geophysical
Research Letters}, \emph{41}(21), 7642--7650.
\url{https://doi.org/10.1002/2014GL061943}

\bibitem[\citeproctext]{ref-Chemke2022}
Chemke, R. (2022). The future poleward shift of {Southern Hemisphere}
summer mid-latitude storm tracks stems from ocean coupling. \emph{Nature
Communications}, \emph{13}(1), 1730.
\url{https://doi.org/10.1038/s41467-022-29392-4}

\bibitem[\citeproctext]{ref-ConsortiumEC-Earth2019}
Consortium (EC-Earth), E.-E. (2019a). {IPCC DDC}: {EC-Earth-Consortium
EC-Earth3} model output prepared for {CMIP6 CMIP} {[}Data set{]}. World
Data Center for Climate (WDCC) at DKRZ.
\url{https://doi.org/10.26050/WDCC/AR6.C6CMEEE3}

\bibitem[\citeproctext]{ref-ConsortiumEC-Earth2019a}
Consortium (EC-Earth), E.-E. (2019b). {IPCC DDC}: {EC-Earth-Consortium
EC-Earth3} model output prepared for {CMIP6 ScenarioMIP} {[}Data set{]}.
World Data Center for Climate (WDCC) at DKRZ.
\url{https://doi.org/10.26050/WDCC/AR6.C6SPEEE3}

\bibitem[\citeproctext]{ref-ConsortiumEC-Earth2019b}
Consortium (EC-Earth), E.-E. (2019c). {IPCC DDC}: {EC-Earth-Consortium
EC-Earth3-Veg} model output prepared for {CMIP6 CMIP} {[}Data set{]}.
World Data Center for Climate (WDCC) at DKRZ.
\url{https://doi.org/10.26050/WDCC/AR6.C6CMEEEVE}

\bibitem[\citeproctext]{ref-ConsortiumEC-Earth2019c}
Consortium (EC-Earth), E.-E. (2019d). {IPCC DDC}: {EC-Earth-Consortium
EC-Earth3-Veg} model output prepared for {CMIP6 ScenarioMIP} {[}Data
set{]}. World Data Center for Climate (WDCC) at DKRZ.
\url{https://doi.org/10.26050/WDCC/AR6.C6SPEEEVE}

\bibitem[\citeproctext]{ref-ConsortiumEC-Earth2020}
Consortium (EC-Earth), E.-E. (2020a). {IPCC DDC}: {EC-Earth-Consortium
EC-Earth-3-CC} model output prepared for {CMIP6 CMIP} {[}Data set{]}.
World Data Center for Climate (WDCC) at DKRZ.
\url{https://doi.org/10.26050/WDCC/AR6.C6CMEEEEC}

\bibitem[\citeproctext]{ref-ConsortiumEC-Earth2020a}
Consortium (EC-Earth), E.-E. (2020b). {IPCC DDC}: {EC-Earth-Consortium
EC-Earth3-Veg-LR} model output prepared for {CMIP6 CMIP} {[}Data set{]}.
World Data Center for Climate (WDCC) at DKRZ.
\url{https://doi.org/10.26050/WDCC/AR6.C6CMEEEVL}

\bibitem[\citeproctext]{ref-ConsortiumEC-Earth2021}
Consortium (EC-Earth), E.-E. (2021). {IPCC DDC}: {EC-Earth-Consortium
EC-Earth3-CC} model output prepared for {CMIP6 ScenarioMIP} {[}Data
set{]}. World Data Center for Climate (WDCC) at DKRZ.
\url{https://doi.org/10.26050/WDCC/AR6.C6SPEEEEC}

\bibitem[\citeproctext]{ref-ConsortiumEC-Earth2023}
Consortium (EC-Earth), E.-E. (2023). {IPCC DDC}: {EC-Earth-Consortium
EC-Earth3-Veg-LR} model output prepared for {CMIP6 ScenarioMIP} {[}Data
set{]}. World Data Center for Climate (WDCC) at DKRZ.
\url{https://doi.org/10.26050/WDCC/AR6.C6SPEEEVL}

\bibitem[\citeproctext]{ref-Corrales2026a}
Corrales, P. (2026a). Derived data for {Using} rapid temperature falls
to estimate future strong cold front frequency in {CMIP6} climate
projections (version v2) {[}Data set{]}. Zenodo.
\url{https://doi.org/10.5281/zenodo.21713532}

\bibitem[\citeproctext]{ref-Corrales2026}
Corrales, P. (2026b, July 31). Code associated with "{Using} rapid
temperature falls to estimate future strong cold front frequency in
{CMIP6} climate projections (version v1.0) {[}software{]}". Zenodo.
\url{https://doi.org/10.5281/zenodo.21713780}

\bibitem[\citeproctext]{ref-Eyring2016}
Eyring, V., Bony, S., Meehl, G. A., Senior, C. A., Stevens, B.,
Stouffer, R. J., \& Taylor, K. E. (2016). Overview of the {Coupled Model
Intercomparison Project Phase} 6 ({CMIP6}) experimental design and
organization. \emph{Geoscientific Model Development}, \emph{9}(5),
1937--1958. \url{https://doi.org/10.5194/gmd-9-1937-2016}

\bibitem[\citeproctext]{ref-Guo2018}
Guo, H., John, J. G., Blanton, C., McHugh, C., Nikonov, S.,
Radhakrishnan, A., Rand, K., Zadeh, N. T., Balaji, V., Durachta, J.,
Dupuis, C., Menzel, R., Robinson, T., Underwood, S., Vahlenkamp, H.,
Bushuk, M., et al. (2018). {IPCC DDC}: {NOAA-GFDL GFDL-CM4} model output
{[}Data set{]}. World Data Center for Climate (WDCC) at DKRZ.
\url{https://doi.org/10.26050/WDCC/AR6.C6CMNGGFC}

\bibitem[\citeproctext]{ref-Guo2018a}
Guo, H., John, J. G., Blanton, C., McHugh, C., Nikonov, S.,
Radhakrishnan, A., Rand, K., Zadeh, N. T., Balaji, V., Durachta, J.,
Dupuis, C., Menzel, R., Robinson, T., Underwood, S., Vahlenkamp, H.,
Dunne, K. A., et al. (2018). {IPCC DDC}: {NOAA-GFDL GFDL-CM4} model
output prepared for {CMIP6 ScenarioMIP} {[}Data set{]}. World Data
Center for Climate (WDCC) at DKRZ.
\url{https://doi.org/10.26050/WDCC/AR6.C6SPNGGFD}

\bibitem[\citeproctext]{ref-Harvey2014}
Harvey, B. J., Shaffrey, L. C., \& Woollings, T. J. (2014).
Equator-to-pole temperature differences and the extra-tropical storm
track responses of the {CMIP5} climate models. \emph{Climate Dynamics},
\emph{43}(5), 1171--1182.
\url{https://doi.org/10.1007/s00382-013-1883-9}

\bibitem[\citeproctext]{ref-Jungclaus2019}
Jungclaus, J., Bittner, M., Wieners, K.-H., Wachsmann, F., Schupfner,
M., Legutke, S., et al. (2019). {IPCC DDC}: {MPI-M MPIESM1}.2-{HR} model
output prepared for {CMIP6 CMIP} {[}Data set{]}. World Data Center for
Climate (WDCC) at DKRZ.
\url{https://doi.org/10.26050/WDCC/AR6.C6CMMXME2}

\bibitem[\citeproctext]{ref-King2024}
King, M. J., Reeder, M. J., \& Jakob, C. (2024). Strong temperature
falls as a cold frontal metric in {Australian} station observations,
reanalyses and climate models. \emph{Quarterly Journal of the Royal
Meteorological Society}, \emph{150}(765), 4788--4805.
\url{https://doi.org/10.1002/qj.4841}

\bibitem[\citeproctext]{ref-Konstali2024}
Konstali, K., Spengler, T., Spensberger, C., \& Sorteberg, A. (2024).
Linking {Future Precipitation Changes} to {Weather Features} in
{CESM2-LE}. \emph{Journal of Geophysical Research: Atmospheres},
\emph{129}(16), e2024JD041190.
\url{https://doi.org/10.1029/2024JD041190}

\bibitem[\citeproctext]{ref-ggside}
Landis, J. (2024). Ggside: Side grammar graphics. Retrieved from
\url{https://github.com/jtlandis/ggside}

\bibitem[\citeproctext]{ref-Lovato2021}
Lovato, T., Peano, D., \& Butenschön, M. (2021a). {IPCC DDC}: {CMCC
CMCC-ESM2} model output prepared for {CMIP6 CMIP} {[}Data set{]}. World
Data Center for Climate (WDCC) at DKRZ.
\url{https://doi.org/10.26050/WDCC/AR6.C6CMCMCMC}

\bibitem[\citeproctext]{ref-Lovato2021a}
Lovato, T., Peano, D., \& Butenschön, M. (2021b). {IPCC DDC}: {CMCC
CMCC-ESM2} model output prepared for {CMIP6 ScenarioMIP} {[}Data set{]}.
World Data Center for Climate (WDCC) at DKRZ.
\url{https://doi.org/10.26050/WDCC/AR6.C6SPCMCMC}

\bibitem[\citeproctext]{ref-Muir2010}
Muir, L. C., \& Reeder, M. J. (2010). Idealized modelling of landfalling
cold fronts. \emph{Quarterly Journal of the Royal Meteorological
Society}, \emph{136}(653), 2147--2161.
\url{https://doi.org/10.1002/qj.692}

\bibitem[\citeproctext]{ref-OGorman2010}
O'Gorman, P. A. (2010). Understanding the varied response of the
extratropical storm tracks to climate change. \emph{Proceedings of the
National Academy of Sciences}, \emph{107}(45), 19176--19180.
\url{https://doi.org/10.1073/pnas.1011547107}

\bibitem[\citeproctext]{ref-ONeill2016}
O'Neill, B. C., Tebaldi, C., van Vuuren, D. P., Eyring, V.,
Friedlingstein, P., Hurtt, G., et al. (2016). The {Scenario Model
Intercomparison Project} ({ScenarioMIP}) for {CMIP6}.
\emph{Geoscientific Model Development}, \emph{9}(9), 3461--3482.
\url{https://doi.org/10.5194/gmd-9-3461-2016}

\bibitem[\citeproctext]{ref-Parfitt2017}
Parfitt, R., Czaja, A., \& Kwon, Y.-O. (2017). The impact of {SST}
resolution change in the {ERA-Interim} reanalysis on wintertime {Gulf
Stream} frontal air-sea interaction. \emph{Geophysical Research
Letters}, \emph{44}(7), 3246--3254.
\url{https://doi.org/10.1002/2017GL073028}

\bibitem[\citeproctext]{ref-patchwork}
Pedersen, T. L. (2024). Patchwork: The composer of plots.

\bibitem[\citeproctext]{ref-Priestley2022}
Priestley, M. D. K., \& Catto, J. L. (2022). Future changes in the
extratropical storm tracks and cyclone intensity, wind speed, and
structure. \emph{Weather and Climate Dynamics}, \emph{3}(1), 337--360.
\url{https://doi.org/10.5194/wcd-3-337-2022}

\bibitem[\citeproctext]{ref-R}
R Core Team. (2023). R: A language and environment for statistical
computing. Vienna, Austria: R Foundation for Statistical Computing.
Retrieved from \url{https://www.R-project.org/}

\bibitem[\citeproctext]{ref-Reeder2015}
Reeder, M. J., Spengler, T., \& Musgrave, R. (2015). Rossby waves,
extreme fronts, and wildfires in southeastern {Australia}.
\emph{Geophysical Research Letters}, \emph{42}(6), 2015--2023.
\url{https://doi.org/10.1002/2015GL063125}

\bibitem[\citeproctext]{ref-Reeder2021}
Reeder, M. J., Spengler, T., \& Spensberger, C. (2021). The {Effect} of
{Sea Surface Temperature Fronts} on {Atmospheric Frontogenesis}.
\emph{Journal of the Atmospheric Sciences}, \emph{78}(6), 1753--1771.
\url{https://doi.org/10.1175/JAS-D-20-0118.1}

\bibitem[\citeproctext]{ref-Schneider2015}
Schneider, T., Bischoff, T., \& Płotka, H. (2015). Physics of {Changes}
in {Synoptic Midlatitude Temperature Variability}. \emph{Journal of
Climate}, \emph{28}(6), 2312--2331.
\url{https://doi.org/10.1175/JCLI-D-14-00632.1}

\bibitem[\citeproctext]{ref-Schupfner2019}
Schupfner, M., Wieners, K.-H., Wachsmann, F., Steger, C., Bittner, M.,
Jungclaus, J., et al. (2019). {IPCC DDC}: {DKRZ MPI-ESM1}.2-{HR} model
output prepared for {CMIP6 ScenarioMIP} {[}Data set{]}. World Data
Center for Climate (WDCC) at DKRZ.
\url{https://doi.org/10.26050/WDCC/AR6.C6SPDKME2}

\bibitem[\citeproctext]{ref-Simmonds2012}
Simmonds, I., Keay, K., \& Bye, J. A. T. (2012). Identification and
{Climatology} of {Southern Hemisphere Mobile Fronts} in a {Modern
Reanalysis}. \emph{Journal of Climate}, \emph{25}(6), 1945--1962.
\url{https://doi.org/10.1175/JCLI-D-11-00100.1}

\bibitem[\citeproctext]{ref-Tamarin2017}
Tamarin, T., \& Kaspi, Y. (2017). The poleward shift of storm tracks
under global warming: {A Lagrangian} perspective. \emph{Geophysical
Research Letters}, \emph{44}(20), 10, 666--10, 674.
\url{https://doi.org/10.1002/2017GL073633}

\bibitem[\citeproctext]{ref-Thomas2019}
Thomas, C. M., \& Schultz, D. M. (2019). Global {Climatologies} of
{Fronts}, {Airmass Boundaries}, and {Airstream Boundaries}: {Why} the
{Definition} of {``{Front}''} {Matters}. \emph{Monthly Weather Review},
\emph{147}(2), 691--717. \url{https://doi.org/10.1175/MWR-D-18-0289.1}

\bibitem[\citeproctext]{ref-Volodin2019}
Volodin, E., Mortikov, E., Gritsun, A., Lykossov, V., Galin, V.,
Diansky, N., et al. (2019a). {INM INM-CM4-8} model output prepared for
{CMIP6 CMIP} historical {[}Data set{]}. Earth System Grid Federation.
\url{https://doi.org/10.22033/ESGF/CMIP6.5069}

\bibitem[\citeproctext]{ref-Volodin2019a}
Volodin, E., Mortikov, E., Gritsun, A., Lykossov, V., Galin, V.,
Diansky, N., et al. (2019b). {IPCC DDC}: {INM INM-CM4-8} model output
prepared for {CMIP6 ScenarioMIP} {[}Data set{]}. World Data Center for
Climate (WDCC) at DKRZ.
\url{https://doi.org/10.26050/WDCC/AR6.C6SPINIC8}

\bibitem[\citeproctext]{ref-Volodin2019b}
Volodin, E., Mortikov, E., Gritsun, A., Lykossov, V., Galin, V.,
Diansky, N., et al. (2019c). {IPCC DDC}: {INM INM-CM5-0} model output
prepared for {CMIP6 CMIP} {[}Data set{]}. World Data Center for Climate
(WDCC) at DKRZ. \url{https://doi.org/10.26050/WDCC/AR6.C6CMINIC0}

\bibitem[\citeproctext]{ref-Volodin2019c}
Volodin, E., Mortikov, E., Gritsun, A., Lykossov, V., Galin, V.,
Diansky, N., et al. (2019d). {IPCC DDC}: {INM INM-CM5-0} model output
prepared for {CMIP6 ScenarioMIP} {[}Data set{]}. World Data Center for
Climate (WDCC) at DKRZ.
\url{https://doi.org/10.26050/WDCC/AR6.C6SPINIC0}

\bibitem[\citeproctext]{ref-ggplot2}
Wickham, H. (2016). ggplot2: Elegant graphics for data analysis.
Springer-Verlag New York. Retrieved from
\url{https://ggplot2.tidyverse.org}

\bibitem[\citeproctext]{ref-cowplot}
Wilke, C. O. (2024). Cowplot: Streamlined plot theme and plot
annotations for 'ggplot2'. Retrieved from
\url{https://wilkelab.org/cowplot/}

\end{CSLReferences}




\end{document}
